% GENERATED FILE. Edit canonical lessons, then run npm run generate:books.
% canonical-source-hash: fnv1a64:b5cc8d8f4027fbac
% canonical-lessons: GU-C461-krpan, GU-C461-udau, GU-C461-kangal, GU-C461-sakhavat, GU-C461-dan

\chapter{Saving and Spending}
\label{ch:gu-saving-and-spending}
\hlchaptermodality{461}

\begin{chapteropening}
\textbf{By the end of this chapter:} \emph{I can say how someone spends money.}

Complete the last lesson of chapter 461: I can say how someone spends money.
\end{chapteropening}

\section[\texorpdfstring{\gu{કૃપણ} (kṛpaṇ)}{કૃપણ (kṛpaṇ)}]{\gu{કૃપણ} (kṛpaṇ) — miserly}

\label{lesson:GU-C461-krpan}



\begin{quote}
Before the new one: say the Gujarati for deceit, then the Gujarati for to cajole.
\end{quote}

\subsection*{You'll want to know: \gu{કૃપણ}}
\textbf{\gu{કૃપણ}} — \emph{kṛpaṇ} — \textquotedblleft{}miserly\textquotedblright{}.

\begin{quote}
\gu{મારા} \gu{કાકા} \gu{ઘણા} \gu{કૃપણ} \gu{છે}.
\end{quote}

\begin{quote}
\emph{mārā kākā ghaṇā kṛpaṇ che.}
\end{quote}

\begin{quote}
My uncle is very miserly.
\end{quote}

\begin{cousinweb}
\textbf{\gu{કૃપણ}} (\emph{kṛpaṇ}) is a tatsama, Sanskrit \emph{kṛpaṇa}, as Monier-Williams records it; its older senses are \textquotedblleft{}pitiable, miserable\textquotedblright{}, and the sense \textquotedblleft{}miserly\textquotedblright{} grew from them.
\end{cousinweb}

\begin{grammarlens}[title={What you've built}]
One more word for describing what you lived and felt.
\end{grammarlens}

\subsection*{Your turn}
\begin{itemize}
\raggedright
  \item \emph{Say it:} \emph{kṛpaṇ}
  \item \emph{Say it:} \emph{kṛpaṇ}, once more
  \item \emph{Say it:} \emph{mārā kākā ghaṇā kṛpaṇ che.} — My uncle is very miserly.
  \item \emph{Recall:} say the Gujarati for deceit, then the Gujarati for to cajole, then say \emph{kṛpaṇ} again
\end{itemize}

\begin{tcolorbox}[breakable,colback=teal!4,colframe=teal!35!black,title={Before you move on}]
What does \gu{કૃપણ} mean? (\textquotedblleft{}Miserly\textquotedblright{}.) Say it once more.
\end{tcolorbox}
\section[\texorpdfstring{\gu{ઉડાઉ} (uḍāu)}{ઉડાઉ (uḍāu)}]{\gu{ઉડાઉ} (uḍāu) — wasteful with money}

\label{lesson:GU-C461-udau}



\begin{quote}
Before the new one: say the Gujarati for to cajole, then the Gujarati for miserly.
\end{quote}

\subsection*{You'll want to know: \gu{ઉડાઉ}}
\textbf{\gu{ઉડાઉ}} — \emph{uḍāu} — \textquotedblleft{}wasteful with money\textquotedblright{}.

\begin{quote}
\gu{ઉડાઉ} \gu{દીકરાએ} \gu{બધા} \gu{પૈસા} \gu{વેડફ્યા}.
\end{quote}

\begin{quote}
\emph{uḍāu dīkrāe badhā paisā veḍaphyā.}
\end{quote}

\begin{quote}
The spendthrift son wasted all the money.
\end{quote}

\begin{cousinweb}
\textbf{\gu{ઉડાઉ}} (\emph{uḍāu}) is probably built on the stem of \textbf{\gu{ઉડાડવું}} (\emph{uḍāḍvũ}), \textquotedblleft{}to make fly\textquotedblright{}, which CDIAL also glosses \textquotedblleft{}to squander\textquotedblright{} (CDIAL 1697), with the adjective ending \emph{-u}: someone who squanders. That verb goes back to Sanskrit \emph{uḍḍīyatē}, \textquotedblleft{}flies up\textquotedblright{}; CDIAL does not list \emph{uḍāu} itself, so the link is probable, not recorded.
\end{cousinweb}

\begin{grammarlens}[title={What you've built}]
One more word for describing what you lived and felt.
\end{grammarlens}

\subsection*{Your turn}
\begin{itemize}
\raggedright
  \item \emph{Say it:} \emph{uḍāu}
  \item \emph{Say it:} \emph{uḍāu}, once more
  \item \emph{Say it:} \emph{uḍāu dīkrāe badhā paisā veḍaphyā.} — The spendthrift son wasted all the money.
  \item \emph{Recall:} say the Gujarati for to cajole, then the Gujarati for miserly, then say \emph{uḍāu} again
\end{itemize}

\begin{tcolorbox}[breakable,colback=teal!4,colframe=teal!35!black,title={Before you move on}]
What does \gu{ઉડાઉ} mean? (\textquotedblleft{}Wasteful with money\textquotedblright{}.) Say it once more.
\end{tcolorbox}
\section[\texorpdfstring{\gu{કંગાળ} (kangāḷ)}{કંગાળ (kangāḷ)}]{\gu{કંગાળ} (kangāḷ) — penniless; destitute}

\label{lesson:GU-C461-kangal}



\begin{quote}
Before the new one: say the Gujarati for miserly, then the Gujarati for wasteful with money.
\end{quote}

\subsection*{You'll want to know: \gu{કંગાળ}}
\textbf{\gu{કંગાળ}} — \emph{kangāḷ} — \textquotedblleft{}penniless; destitute\textquotedblright{}.

\begin{quote}
\gu{ધંધામાં} \gu{ખોટ} \gu{ગઈ} \gu{અને} \gu{તે} \gu{કંગાળ} \gu{થઈ} \gu{ગયો}.
\end{quote}

\begin{quote}
\emph{dhandhāmā̃ khoṭ gaī ane te kangāḷ thaī gayo.}
\end{quote}

\begin{quote}
His business made a loss and he became penniless.
\end{quote}

\begin{cousinweb}
\textbf{\gu{કંગાળ}} (\emph{kangāḷ}) goes back to a reconstructed \emph{*kaṅkāla-}, \textquotedblleft{}poor, miserable\textquotedblright{}; CDIAL lists the Gujarati word, \textquotedblleft{}poor\textquotedblright{} (CDIAL 2604).
\end{cousinweb}

\begin{grammarlens}[title={What you've built}]
One more word for describing what you lived and felt.
\end{grammarlens}

\subsection*{Your turn}
\begin{itemize}
\raggedright
  \item \emph{Say it:} \emph{kangāḷ}
  \item \emph{Say it:} \emph{kangāḷ}, once more
  \item \emph{Say it:} \emph{dhandhāmā̃ khoṭ gaī ane te kangāḷ thaī gayo.} — His business made a loss and he became penniless.
  \item \emph{Recall:} say the Gujarati for miserly, then the Gujarati for wasteful with money, then say \emph{kangāḷ} again
\end{itemize}

\begin{tcolorbox}[breakable,colback=teal!4,colframe=teal!35!black,title={Before you move on}]
What does \gu{કંગાળ} mean? (\textquotedblleft{}Penniless; destitute\textquotedblright{}.) Say it once more.
\end{tcolorbox}
\section[\texorpdfstring{\gu{સખાવત} (sakhāvat)}{સખાવત (sakhāvat)}]{\gu{સખાવત} (sakhāvat) — charity; open-handed giving}

\label{lesson:GU-C461-sakhavat}



\begin{quote}
Before the new one: say the Gujarati for wasteful with money, then the Gujarati for penniless.
\end{quote}

\subsection*{You'll want to know: \gu{સખાવત}}
\textbf{\gu{સખાવત}} — \emph{sakhāvat} — \textquotedblleft{}charity; open-handed giving\textquotedblright{}.

\begin{quote}
\gu{તે} \gu{વેપારીએ} \gu{ગામ} \gu{માટે} \gu{ઘણી} \gu{સખાવત} \gu{કરી}.
\end{quote}

\begin{quote}
\emph{te vepārīe gām māṭe ghaṇī sakhāvat karī.}
\end{quote}

\begin{quote}
That merchant gave generously to the village.
\end{quote}

A feminine noun.

\begin{cousinweb}
\textbf{\gu{સખાવત}} (\emph{sakhāvat}) is usually traced to Arabic \emph{saxāwa}, \textquotedblleft{}generosity\textquotedblright{}, which reached Gujarati through Persian.
\end{cousinweb}

\begin{grammarlens}[title={What you've built}]
One more word for describing what you lived and felt.
\end{grammarlens}

\subsection*{Your turn}
\begin{itemize}
\raggedright
  \item \emph{Say it:} \emph{sakhāvat}
  \item \emph{Say it:} \emph{sakhāvat}, once more
  \item \emph{Say it:} \emph{te vepārīe gām māṭe ghaṇī sakhāvat karī.} — That merchant gave generously to the village.
  \item \emph{Recall:} say the Gujarati for wasteful with money, then the Gujarati for penniless, then say \emph{sakhāvat} again
\end{itemize}

\begin{tcolorbox}[breakable,colback=teal!4,colframe=teal!35!black,title={Before you move on}]
What does \gu{સખાવત} mean? (\textquotedblleft{}Charity; open-handed giving\textquotedblright{}.) Say it once more.
\end{tcolorbox}
\section[\texorpdfstring{\gu{દાન} (dān)}{દાન (dān)}]{\gu{દાન} (dān) — a donation; a gift to charity}

\label{lesson:GU-C461-dan}



\begin{quote}
Before the new one: say the Gujarati for penniless, then the Gujarati for charity.
\end{quote}

\subsection*{You'll want to know: \gu{દાન}}
\textbf{\gu{દાન}} — \emph{dān} — \textquotedblleft{}a donation; a gift to charity\textquotedblright{}.

\begin{quote}
\gu{દાદાએ} \gu{મંદિરમાં} \gu{દાન} \gu{આપ્યું}.
\end{quote}

\begin{quote}
\emph{dādāe mandirmā̃ dān āpyũ.}
\end{quote}

\begin{quote}
Grandfather made a donation at the temple.
\end{quote}

A neuter noun.

\begin{cousinweb}
\textbf{\gu{દાન}} (\emph{dān}) is a tatsama, Sanskrit \emph{dāna}, \textquotedblleft{}the act of giving\textquotedblright{}, as Monier-Williams records it.
\end{cousinweb}

\begin{grammarlens}[title={What you've built}]
One more word for describing what you lived and felt.
\end{grammarlens}

\subsection*{Your turn}
\begin{itemize}
\raggedright
  \item \emph{Say it:} \emph{dān}
  \item \emph{Say it:} \emph{dān}, once more
  \item \emph{Say it:} \emph{dādāe mandirmā̃ dān āpyũ.} — Grandfather made a donation at the temple.
  \item \emph{Recall:} say the Gujarati for penniless, then the Gujarati for charity, then say \emph{dān} again
\end{itemize}

\begin{tcolorbox}[breakable,colback=teal!4,colframe=teal!35!black,title={Before you move on}]
What does \gu{દાન} mean? (\textquotedblleft{}A donation; a gift to charity\textquotedblright{}.) Say it once more.
\end{tcolorbox}
